% Lecture 8: Midterm review, mock exam, and walkthrough setup.
% Compile from the repository root with:
%   pdflatex -interaction=nonstopmode -halt-on-error -output-directory=output/pdf Lecture8/lecture8.tex
\pdfminorversion=7
\documentclass[aspectratio=169]{beamer}

\usepackage{amsmath,amssymb,mathtools}
\usepackage{array,booktabs}
\usepackage{xcolor}
\usepackage{iftex}
\ifPDFTeX
  \usepackage[UTF8,heading=false]{ctex}
\else
  \usepackage{fontspec}
  \usepackage{xeCJK}
  \setCJKmainfont{Microsoft YaHei}
\fi

% Allow compilation both from the repository root and from Lecture8/.
% The custom Beamer theme is stored at the repository root.
\makeatletter
\def\input@path{{./}{../}{Lecture8/}}
\makeatother

\mode<presentation>{
  \usetheme{OSUCSE169}
  \setbeamercovered{invisible}
}
\setbeamertemplate{navigation symbols}{}
\setbeamertemplate{footline}[frame number]
\setbeamercolor{headline}{fg=black,bg=black}
\setbeamercolor{section in head/foot}{bg=gray}
\setbeamertemplate{headline}{%
  \begin{beamercolorbox}[ht=2.25ex,dp=3.75ex]{section in head/foot}
    \insertnavigation{\paperwidth}
  \end{beamercolorbox}}

\definecolor{lecturegreen}{RGB}{48,105,21}
\definecolor{lectureblue}{RGB}{0,76,170}
\definecolor{lectureorange}{RGB}{202,105,0}
\definecolor{softgray}{RGB}{245,245,245}

\newcommand{\R}{\mathbb{R}}
\newcommand{\vect}[1]{\mathbf{#1}}
\newcommand{\Span}{\operatorname{Span}}
\newcommand{\Col}{\operatorname{Col}}
\newcommand{\Nul}{\operatorname{Nul}}
\newcommand{\rank}{\operatorname{rank}}
\newcommand{\nullity}{\operatorname{nullity}}
\newcommand{\keybox}[2]{%
  \begin{beamercolorbox}[sep=1.0ex,wd=\linewidth]{block body}
    \textbf{#1}\par #2
  \end{beamercolorbox}}
\newcommand{\sectioncard}[2]{%
  \begin{frame}[plain]
    \vfill\centering
    {\Huge\color{lectureorange}\bfseries #1}\par\bigskip
    {\LARGE\color{lecturegreen}\bfseries #2}
    \vfill
  \end{frame}}

\title{MATH 100: Linear Algebra}
\subtitle{Lecture 8 - Midterm Review}
\author{Yuqing Yang}
\institute{Macau University of Science and Technology}
\date{10/21/2026}

\begin{document}

\begin{frame}[plain]
  \titlepage
\end{frame}

\begin{frame}{Today's session}
\begin{enumerate}
  \item Review the main ideas from Lectures 2--7.
  \item Complete a mock midterm under exam conditions.
  \item Work through the mock exam and compare solution methods.
\end{enumerate}
\medskip
\begin{block}{Lecture 1 reminder}
The midterm is closed-book. Use the review to identify topics that still require practice.
\end{block}
\end{frame}

\section{Systems and transformations}
\sectioncard{Part I}{Systems, vectors, and transformations}

\begin{frame}{The central equation}
\centering
{\LARGE
\[
A\vect x=\vect b
\]}
\vspace{-1ex}
\begin{columns}[T]
\column{.48\textwidth}
\textbf{Existence}
\begin{itemize}
  \item Is the system consistent?
  \item Is $\vect b$ in the span of the columns?
  \item Is $\vect b$ in $\Col A$?
\end{itemize}
\column{.48\textwidth}
\textbf{Uniqueness}
\begin{itemize}
  \item Are there free variables?
  \item Is $\Nul A=\{\vect0\}$?
  \item Are the columns independent?
\end{itemize}
\end{columns}
\end{frame}

\begin{frame}{Linear combinations and span}
For $\vect v_1,\ldots,\vect v_p\in\R^n$,
\[
\Span\{\vect v_1,\ldots,\vect v_p\}
=\{c_1\vect v_1+\cdots+c_p\vect v_p:c_i\in\R\}.
\]
If $A=[\vect a_1\ \cdots\ \vect a_n]$, then
\[
\vect b\in\Span\{\vect a_1,\ldots,\vect a_n\}
\iff A\vect x=\vect b\text{ is consistent}.
\]
\begin{block}{Computational test}
Row reduce $[A\mid\vect b]$. A row $[0\ \cdots\ 0\mid d]$ with $d\ne0$ proves that $\vect b$ is not in the span.
\end{block}
\end{frame}

\begin{frame}{Three forms of one system}
The following descriptions have the same solution set:
\begin{enumerate}
  \item a system of linear equations;
  \item a vector equation,
  \[
  x_1\vect a_1+\cdots+x_n\vect a_n=\vect b;
  \]
  \item a matrix equation,
  \[
  A\vect x=\vect b.
  \]
\end{enumerate}
For $A\in\R^{m\times n}$, the input is in $\R^n$ and the output is in $\R^m$.
\end{frame}

\begin{frame}{Pivots, free variables, and solution types}
\renewcommand{\arraystretch}{1.35}
\begin{center}
\begin{tabular}{>{\raggedright\arraybackslash}p{.34\linewidth}>{\raggedright\arraybackslash}p{.51\linewidth}}
\toprule
Row-reduction result & Conclusion \\
\midrule
$[0\ \cdots\ 0\mid d]$, $d\ne0$ & No solution \\
Consistent, no free variables & One solution \\
Consistent, at least one free variable & Infinitely many solutions \\
\bottomrule
\end{tabular}
\end{center}
\medskip
For $A\vect x=\vect0$:
\[
\text{nontrivial solutions exist}
\iff \text{at least one free variable exists}.
\]
\end{frame}

\begin{frame}{Parametric vector form}
For a homogeneous system,
\[
\vect x=s\vect v_1+t\vect v_2.
\]
For a consistent nonhomogeneous system,
\[
\vect x=\vect p+s\vect v_1+t\vect v_2,
\]
where $\vect p$ is one particular solution and the parameter vectors solve $A\vect x=\vect0$.
\begin{block}{Workflow}
Identify free variables, introduce parameters, solve for pivot variables, then separate the constant and parameter vectors.
\end{block}
\end{frame}

\begin{frame}{Linear independence}
$\{\vect v_1,\ldots,\vect v_p\}$ is linearly independent when
\[
x_1\vect v_1+\cdots+x_p\vect v_p=\vect0
\]
has only the trivial solution.
\begin{itemize}
  \item Put the vectors into the columns of $A$ and solve $A\vect x=\vect0$.
  \item A set containing $\vect0$ is dependent.
  \item Two vectors are dependent exactly when one is a scalar multiple of the other.
  \item More than $n$ vectors in $\R^n$ are dependent.
  \item A dependent set contains a vector that is a linear combination of the others.
\end{itemize}
\end{frame}

\begin{frame}{Linear transformations}
A transformation $T$ is linear when
\[
T(\vect u+\vect v)=T(\vect u)+T(\vect v),
\qquad
T(c\vect u)=cT(\vect u).
\]
Therefore,
\[
T(c_1\vect v_1+\cdots+c_p\vect v_p)
=c_1T(\vect v_1)+\cdots+c_pT(\vect v_p).
\]
Every matrix transformation $T(\vect x)=A\vect x$ is linear.
\begin{alertblock}{Quick rejection test}
If $T(\vect0)\ne\vect0$, then $T$ is not linear.
\end{alertblock}
\end{frame}

\begin{frame}{The standard matrix}
Every linear transformation $T:\R^n\to\R^m$ has a unique matrix
\[
A=[T(\vect e_1)\ \cdots\ T(\vect e_n)]
\]
such that $T(\vect x)=A\vect x$.
\begin{block}{Composition order}
If $A$ acts first and $B$ acts second, then the composition has matrix $BA$.
\end{block}
The standard matrices studied in class describe dilations, contractions, shears, projections, and reflections.
\end{frame}

\begin{frame}{Onto and one-to-one}
For $T(\vect x)=A\vect x$, with $A\in\R^{m\times n}$:
\begin{columns}[T]
\column{.48\textwidth}
\textbf{Onto $\R^m$}
\begin{itemize}
  \item $A\vect x=\vect b$ is consistent for every $\vect b\in\R^m$.
  \item $A$ has a pivot in every row.
\end{itemize}
\column{.48\textwidth}
\textbf{One-to-one}
\begin{itemize}
  \item $A\vect x=\vect0$ has only the trivial solution.
  \item $A$ has a pivot in every column.
\end{itemize}
\end{columns}
Also,
\[
\ker T=\Nul A,
\qquad
\operatorname{range}T=\Col A.
\]
\end{frame}

\section{Matrix operations}
\sectioncard{Part II}{Matrix operations and invertibility}

\begin{frame}{Matrix multiplication}
The product is defined when the inner dimensions match:
\[
A_{m\times r}B_{r\times n}=C_{m\times n}.
\]
Three useful views:
\begin{itemize}
  \item each entry is a row-column dot product;
  \item each column of $AB$ is $A$ times the corresponding column of $B$;
  \item in $AB\vect x$, the matrix $B$ acts first.
\end{itemize}
\[
A(BC)=(AB)C,
\quad
A(B+C)=AB+AC.
\]
In general, $AB\ne BA$, and $AB=0$ does not force $A=0$ or $B=0$.
\end{frame}

\begin{frame}{Transpose and special matrices}
\[
(A^T)_{ij}=A_{ji}.
\]
\begin{align*}
(A^T)^T&=A,&
(A+B)^T&=A^T+B^T,\\
(cA)^T&=cA^T,&
(AB)^T&=B^TA^T.
\end{align*}
\begin{itemize}
  \item $A$ is symmetric when $A^T=A$.
  \item $A$ is skew-symmetric when $A^T=-A$.
  \item The diagonal entries of a skew-symmetric matrix are zero.
\end{itemize}
\end{frame}

\begin{frame}{Inverse matrices}
A square matrix is invertible when
\[
AA^{-1}=I=A^{-1}A.
\]
\begin{align*}
(A^{-1})^{-1}&=A,\\
(AB)^{-1}&=B^{-1}A^{-1},\\
(A^T)^{-1}&=(A^{-1})^T.
\end{align*}
Row-reduction algorithm:
\[
[A\mid I]\longrightarrow[I\mid A^{-1}].
\]
If the left block cannot become $I$, then $A$ is singular.
\end{frame}

\begin{frame}{Matrix equations}
When the required inverses exist,
\begin{align*}
AX=B &\Longrightarrow X=A^{-1}B,\\
XA=B &\Longrightarrow X=BA^{-1},\\
AXB=C &\Longrightarrow X=A^{-1}CB^{-1}.
\end{align*}
\begin{alertblock}{Order matters}
Multiply by an inverse on the side where the matching factor appears. Do not commute matrix factors.
\end{alertblock}
\end{frame}

\begin{frame}[shrink=7]{The Invertible Matrix Theorem}
For a square $n\times n$ matrix $A$, the following are equivalent:
\begin{columns}[T]
\column{.49\textwidth}
\begin{itemize}
  \item $A$ is invertible.
  \item $A\sim I_n$.
  \item $A$ has $n$ pivots.
  \item $A\vect x=\vect0$ has only the trivial solution.
  \item The columns are independent.
  \item The columns span $\R^n$.
\end{itemize}
\column{.49\textwidth}
\begin{itemize}
  \item $A\vect x=\vect b$ has one solution for every $\vect b$.
  \item $\vect x\mapsto A\vect x$ is one-to-one and onto.
  \item $A^T$ is invertible.
  \item $\rank A=n$.
  \item $\Nul A=\{\vect0\}$.
  \item $\det A\ne0$.
\end{itemize}
\end{columns}
\end{frame}

\begin{frame}{Partitioned matrices}
Block multiplication follows the usual row-column rule:
\[
\begin{bmatrix}A_{11}&A_{12}\\A_{21}&A_{22}\end{bmatrix}
\begin{bmatrix}B_{11}&B_{12}\\B_{21}&B_{22}\end{bmatrix}
=
\begin{bmatrix}
A_{11}B_{11}+A_{12}B_{21}&A_{11}B_{12}+A_{12}B_{22}\\
A_{21}B_{11}+A_{22}B_{21}&A_{21}B_{12}+A_{22}B_{22}
\end{bmatrix}.
\]
Every block product must have compatible dimensions.
\[
AB=\sum_{k=1}^{n}\operatorname{col}_k(A)\operatorname{row}_k(B).
\]
\end{frame}

\begin{frame}{LU factorization}
Write
\[
A=LU,
\]
where $L$ is unit lower triangular and $U$ is echelon or upper triangular.
\begin{enumerate}
  \item Reduce $A$ to $U$ using row replacements.
  \item Record the elimination multipliers in $L$.
  \item Verify $LU=A$.
\end{enumerate}
To solve $A\vect x=\vect b$:
\[
L\vect y=\vect b
\quad\text{then}\quad
U\vect x=\vect y.
\]
The factorization can be reused for many right-hand sides.
\end{frame}

\section{Determinants}
\sectioncard{Part III}{Determinants and their applications}

\begin{frame}{Determinant computation}
For a $2\times2$ matrix,
\[
\det\begin{bmatrix}a&b\\c&d\end{bmatrix}=ad-bc.
\]
For larger matrices,
\[
C_{ij}=(-1)^{i+j}\det A_{ij},
\qquad
\det A=\sum_j a_{ij}C_{ij}.
\]
\begin{itemize}
  \item Expand along a row or column with many zeros.
  \item If $A$ is triangular, multiply its diagonal entries.
  \item Sarrus' rule applies only to $3\times3$ matrices.
\end{itemize}
\end{frame}

\begin{frame}{Row operations and determinants}
\renewcommand{\arraystretch}{1.4}
\begin{center}
\begin{tabular}{>{\raggedright\arraybackslash}p{.46\linewidth}>{\raggedright\arraybackslash}p{.34\linewidth}}
\toprule
Row operation & Effect on determinant \\
\midrule
$R_i\leftrightarrow R_j$ & Multiply by $-1$ \\
$R_i\leftarrow cR_i$ & Multiply by $c$ \\
$R_i\leftarrow R_i+cR_j$ & No change \\
\bottomrule
\end{tabular}
\end{center}
The same rules hold for column operations. A zero row, two equal rows, or proportional rows force the determinant to be zero.
\end{frame}

\begin{frame}{Determinant identities}
For square $n\times n$ matrices,
\begin{align*}
\det(AB)&=\det A\det B,\\
\det(A^T)&=\det A,\\
\det(cA)&=c^n\det A,\\
\det(A^k)&=(\det A)^k.
\end{align*}
If $A$ is invertible,
\[
\det(A^{-1})=\frac{1}{\det A}.
\]
\begin{alertblock}{False in general}
$\det(A+B)=\det A+\det B$.
\end{alertblock}
\end{frame}

\begin{frame}{Determinants and system classification}
For square $A$,
\[
A\text{ is invertible}\iff\det A\ne0.
\]
Hence $\det A\ne0$ gives one solution for every $\vect b$.
\medskip

If $\det A=0$, compare ranks:
\begin{align*}
\rank A=\rank[A\mid\vect b]<n
&\Longrightarrow \text{infinitely many solutions},\\
\rank A<\rank[A\mid\vect b]
&\Longrightarrow \text{no solution}.
\end{align*}
The equation $\det A=0$ does not by itself prove inconsistency.
\end{frame}

\begin{frame}{Cramer's rule}
If $A$ is invertible, then the solution of $A\vect x=\vect b$ satisfies
\[
x_i=\frac{\det A_i(\vect b)}{\det A},
\]
where $A_i(\vect b)$ replaces column $i$ of $A$ with $\vect b$.
\begin{itemize}
  \item Confirm that $A$ is square and $\det A\ne0$.
  \item Replace a column, not a row.
  \item Use elimination for larger or general systems.
\end{itemize}
\end{frame}

\section{Subspaces and rank}
\sectioncard{Part IV}{Subspaces, bases, dimension, and rank}

\begin{frame}{Subspaces}
A subset $H\subseteq\R^n$ is a subspace when
\begin{enumerate}
  \item $\vect0\in H$;
  \item $\vect u+\vect v\in H$ whenever $\vect u,\vect v\in H$;
  \item $c\vect u\in H$ whenever $\vect u\in H$.
\end{enumerate}
Examples include every span, $\Col A$, $\Nul A$, $\R^n$, and $\{\vect0\}$.
\begin{alertblock}{Origin test}
A line or plane that does not pass through the origin is not a subspace.
\end{alertblock}
\end{frame}

\begin{frame}{Column space and null space}
For $A\in\R^{m\times n}$,
\[
\Col A=\Span\{\vect a_1,\ldots,\vect a_n\}\subseteq\R^m,
\]
\[
\Nul A=\{\vect x\in\R^n:A\vect x=\vect0\}\subseteq\R^n.
\]
Membership tests:
\begin{align*}
\vect b\in\Col A
&\iff A\vect x=\vect b\text{ is consistent},\\
\vect x\in\Nul A
&\iff A\vect x=\vect0.
\end{align*}
\end{frame}

\begin{frame}{Bases for column and null spaces}
\textbf{Basis for $\Col A$}
\begin{enumerate}
  \item Row reduce $A$ and locate the pivot columns.
  \item Take the corresponding columns from the original matrix.
\end{enumerate}
\medskip
\textbf{Basis for $\Nul A$}
\begin{enumerate}
  \item Solve $A\vect x=\vect0$.
  \item Write the solution in parametric vector form.
  \item Use one coefficient vector for each free parameter.
\end{enumerate}
\begin{alertblock}{Column-space warning}
Do not use the columns of the reduced matrix as a basis for the original column space.
\end{alertblock}
\end{frame}

\begin{frame}{Coordinates and dimension}
If $\beta=\{\vect b_1,\ldots,\vect b_p\}$ is a basis for $H$, then every $\vect x\in H$ has a unique representation
\[
\vect x=c_1\vect b_1+\cdots+c_p\vect b_p.
\]
Its coordinate vector is
\[
[\vect x]_\beta=\begin{bmatrix}c_1\\\vdots\\c_p\end{bmatrix}.
\]
The dimension of $H$ is the number of vectors in any basis:
\[
\dim H=p.
\]
\end{frame}

\begin{frame}{Rank and nullity}
\[
\rank A=\dim(\Col A)=\text{number of pivot columns},
\]
\[
\nullity A=\dim(\Nul A)=\text{number of free variables}.
\]
If $A$ has $n$ columns, then
\[
\boxed{\rank A+\nullity A=n}.
\]
For an $n\times n$ matrix,
\[
A\text{ invertible}
\iff \rank A=n
\iff \nullity A=0.
\]
\end{frame}

\begin{frame}{Choosing a method}
\small
\renewcommand{\arraystretch}{1.27}
\begin{center}
\begin{tabular}{>{\raggedright\arraybackslash}p{.46\linewidth}>{\raggedright\arraybackslash}p{.42\linewidth}}
\toprule
Question & First move \\
\midrule
$\vect b$ in a span or column space? & Reduce $[A\mid\vect b]$ \\
Columns independent? & Solve $A\vect x=\vect0$ \\
Onto or one-to-one? & Check row or column pivots \\
Basis for $\Col A$? & Original pivot columns \\
Basis for $\Nul A$? & Parametric homogeneous solution \\
Square matrix invertible? & Easiest equivalent condition \\
Large determinant? & Tracked elimination \\
Repeated systems with one $A$? & LU factorization \\
\bottomrule
\end{tabular}
\end{center}
\end{frame}

\begin{frame}{Common mistakes}
\begin{itemize}
  \item Keeping the factor order when transposing or inverting a product.
  \item Treating matrix multiplication as commutative.
  \item Assuming $\det(A+B)=\det A+\det B$.
  \item Forgetting $\det(cA)=c^n\det A$.
  \item Concluding that $\det A=0$ means no solution.
  \item Taking column-space basis vectors from the reduced matrix.
  \item Confusing a pivot in every row with a pivot in every column.
  \item Using Cramer's rule for a singular or rectangular matrix.
\end{itemize}
\end{frame}

\begin{frame}{Review checkpoint}
Can you answer these questions without notes?
\begin{enumerate}
  \item How do pivots classify a linear system?
  \item How do you test span, independence, onto, and one-to-one?
  \item How do you find inverse, column-space basis, and null-space basis?
  \item How do row operations change a determinant?
  \item Which statements are equivalent to invertibility?
  \item How do rank and nullity divide the variables?
\end{enumerate}
\end{frame}

\section{Mock exam}
\sectioncard{Part V}{Mock midterm}

\begin{frame}{Mock exam format}
\begin{center}
\renewcommand{\arraystretch}{1.4}
\begin{tabular}{lcc}
\toprule
Section & Questions & Points \\
\midrule
Multiple choice & 10 & 20 \\
Brief answers & 5 & 10 \\
Exercises & 7 & 70 \\
\midrule
Total & 22 & 100 \\
\bottomrule
\end{tabular}
\end{center}
\medskip
Show enough work to make your reasoning clear. The walkthrough begins after the exam session ends.
\end{frame}

\begin{frame}{Walkthrough focus}
For each problem, we will compare
\begin{itemize}
  \item the mathematical idea being tested;
  \item the shortest reliable method;
  \item the steps that earn partial credit;
  \item the most likely incorrect approach.
\end{itemize}
\medskip
Bring questions from the mock exam into the walkthrough.
\end{frame}

\end{document}
